\section{Instantiation}\label{sec:offline-instantiation}

We now instantiate the generic construction of \Cref{sec:offline-construction}
with concrete cryptographic building blocks. We use BBS+ signatures to build the
anonymous credential that binds a user's pseudonym to their secret key. We use
Pointcheval-Sanders signatures, extended to support blind signing over committed
inputs, to let the central bank sign tokens at withdrawal without learning their
content. Both are pairing based multi-message signature schemes.  This lets us
build efficient sigma protocols directly over their verification equations,
rather than relying on general purpose zero knowledge proofs. The same pair of
schemes underlies the transaction system of~\Cref{chap:txid}, so the two
chapters instantiate the same infrastructure.

\paragraph{Notation.} We use the notation of the preceding chapters. Groups
$\group{1}$, $\group{2}$, and $\group{T}$ have prime order $p$. They are
equipped with an efficiently computable pairing $\pairing \colon \group{1}
\times \group{2} \rightarrow \group{T}$, and we work in the asymmetric setting.
Elements of $\group{1}$ are written using lower case letters, elements of
$\group{2}$ are distinguished using a tilde, and elements of $\Zp$ are written
using lower case Greek letters. We use generators $\gen \in \group{1}$ and
$\ggen \in \group{2}$, and hash functions $\hash_p$ into $\Zp$ and
$\hash_{\group{1}}$ into $\group{1}$. Every generator used as part of a
commitment key is obtained as $\hash_{\group{1}}$ of a fixed string, so that no
participant knows the discrete logarithm relations among them.

\subsection{Joint proofs of knowledge}\label{sec:joint-pok-def}

Each payment carries a proof that the payer holds a credential on the identifier
committed in their pseudonym, and knows the secret key that credential is bound
to. The secret key never leaves the secure element, and the secure element does
not have the compute to produce the whole proof. We therefore need a proof
system in which two provers cooperate to produce a single proof for a relation
whose witness is split between them, with the weaker prover holding only a small
part of it.

Such a primitive, which we call a joint proof of knowledge, is given
in~\cite{USENIX:HesSinSor23}. We recall it here in the syntax used throughout
this thesis. It is a proof system for a relation $\rel$ over instance $\inst$
and witness $\witness = (\witness_0, \witness_1)$, where $\witness_0$ is held by
prover $\algo{Prv}_0$ and $\witness_1$ by prover $\algo{Prv}_1$.

\begin{primitivebox}{Joint proof of knowledge}{jpok}

  \begin{description}

    \item[$\crs \sample \algo{JGen}(\secparam)$ :] The setup algorithm, on input
security parameter $\secparam$, outputs a common reference string $\crs$.

    \item[$\Pi \sample \algo{JProve}_0(\crs, \inst, \witness_0) \leftrightarrow
\algo{JProve}_1(\crs, \inst, \witness_1)$ :] The proving protocol, run between
$\algo{Prv}_0$ on input witness share $\witness_0$ and $\algo{Prv}_1$ on input
witness share $\witness_1$, with $\algo{Prv}_1$ outputting proof $\Pi$.

    \item[$\{0,1\} \leftarrow \algo{JVerify}(\crs, \inst, \Pi)$ :] The
verification algorithm, on input $\crs$, instance $\inst$, and proof $\Pi$,
returns 0 or 1.

  \end{description}

\end{primitivebox}

A joint proof of knowledge should satisfy the following three properties, which
are those established for the primitive in~\cite{USENIX:HesSinSor23}.

\paragraph{Completeness.} If both provers follow the protocol on a witness
$\witness = (\witness_0, \witness_1)$ with $\rel(\inst, \witness) = 1$, then
$\algo{JVerify}(\crs, \inst, \Pi) = 1$.

\paragraph{Soundness.} There is an efficient extractor that, given a proof
accepted by $\algo{JVerify}$, outputs a witness $\witness$ for $\inst$. This
holds when $\algo{Prv}_0$ and $\algo{Prv}_1$ are both corrupt and colluding,
which is the case we need: a user has full control of their mobile device, and a
user whose secure element is also compromised controls both provers.

\paragraph{Zero knowledge.} There is a simulator that produces the verifier's
view of the proving protocol without access to either witness share.

We rely on this primitive in three ways that go slightly beyond the statement of
the result we cite, and record them here.

First, zero knowledge is proven in~\cite{USENIX:HesSinSor23} against a
semi-honest verifier, on the grounds that in their protocol the verifier relays
messages between the two provers and can therefore be caught deviating by the
honest prover. Our payee may be arbitrarily malicious, but is not interposed
between the payer's secure element and their device: the two provers communicate
directly, and the payee contributes only the instance component $\com_{i+1}$
before receiving the finished proof. The restriction is therefore not binding in
our setting.

Second, $\algo{Prv}_1$ learns $\witness_0$ only through a commitment to it,
which it receives together with that commitment's randomness. The commitment is
consequently not hiding, and what protects $\witness_0$ from the co-prover is
the hardness of computing discrete logarithms rather than any property of the
commitment scheme. This is what lets the mobile device carry out the bulk of the
proof without ever holding the user's secret key, which is the point of the
split: the device is under the user's full control, while the secure element is
not.

Third, the simulator of~\Cref{sec:offline-security} extracts from proofs
produced by corrupt users. Most of these extractions serve to derive a
contradiction, and are performed once before the reduction terminates, so
extraction by rewinding suffices. The exception is the withdrawal proof, from
which the simulator must extract in order to continue the simulation, on
instances a corrupt user chooses adaptively and without bound. There we require
the extractor to be online. This mirrors the treatment
in~\cite{USENIX:HesSinSor23}, where online extractability is assumed only for
the proof accompanying blind issuance. The withdrawal proof is produced by the
mobile device and not by the secure element, so this requirement does not bear
on what the constrained party computes.

\subsection{Instantiating the joint proof of knowledge}\label{sec:joint-proof}

We instantiate the joint proof of knowledge over BBS+ signatures. A BBS+
signature lets a signer sign several messages at once. It is existentially
unforgeable under the $q$-strong Diffie-Hellman ($q$-SDH) assumption. We use it
exactly as described in~\Cref{sec:txid-instantiation}, and do not repeat the
algorithms here. We first recall the standard proof of knowledge of a BBS+
signature, then split it between two provers.

\paragraph{Proof of knowledge of a BBS+ signature.}

Let $\phi = (\varepsilon, s, e) \gets \BBSsign(\sk_\BBS, \msg)$. By construction
$e = (\hgen \vec{q}^{\msg\|s})^{1/(\theta+\varepsilon)}$ for $\sk_\BBS =
\theta$. Setting $\inst = \pk_\BBS = (\ggen, \tilde{t}, \vec{q}, \hgen)$ and
$\witness = (\msg, \phi)$, a proof of knowledge of a BBS+ signature shows that
$(\inst, \witness)$ satisfies the equation as follows:
%
\begin{equation*}
%
  \pairing(e, \ggen^\varepsilon \tilde{t}) = \pairing(\hgen \vec{q}^{\msg\|s},
\ggen).
  %
\end{equation*}
%
A prover can instead select $\iota, \lambda$ at random and compute $u =
e^{-\iota/\lambda}$, $w = (\hgen \vec{q}^{\msg\|s})^{-1/\lambda}$, and $x =
u^\varepsilon w^\iota$. It then proves that $\inst = (\pk_\BBS, u, w, x)$ and
$\witness = (\varepsilon, s, \iota, \lambda, \msg)$ satisfy the relation as
follows, which is equivalent to the verification equation
above~\cite{EC:TesZhu23b}:
%
\begin{align}
%
  \pairing(u, \tilde{t}) \pairing(x, \ggen) = 1 \ \wedge \ x = u^\varepsilon
w^\iota \ \wedge \ \hgen^{-1} = w^\lambda \vec{q}^{\msg\|s}.  \label{eq:BBSrel}
%
\end{align}
%
A prover $\algo{Prv}$ for this relation follows a sigma protocol made non
interactive via the Fiat-Shamir heuristic~\cite{C:FiaSha86}.

Concretely, $\algo{Prv}$ selects $\alpha, \beta, \delta, \gamma_0, \ldots,
\gamma_{\ell-1}, \gamma \in \Zp$ and computes the following:
%
\begin{align*}
%
  a &= u^\alpha w^\beta \ ; \ b = w^\delta q_\ell^\gamma \prod_{i=0}^{\ell-1}
q_i^{\gamma_i} \ ; \ \algo{chal} = \hash_p(\inst, a, b) \\
  %
  \eta &= \alpha + \algo{chal} \varepsilon \ ; \ \nu = \beta + \algo{chal} \iota
\ ; \ \xi = \delta + \algo{chal} \lambda \ ; \\
  %
  \pi_i &= \gamma_i + \algo{chal} \msg_i, \ \forall i \in [\ell] \ ; \ \pi =
\gamma + \algo{chal} s.
  %
\end{align*}
%
It then outputs $\Pi = (a, b, \eta, \nu, \xi, \pi_0, \ldots, \pi_{\ell-1},
\pi)$.  A verifier $\algo{V}$ first checks that $\pairing(u, \tilde{t})
\pairing(x, \ggen) = 1$, computes $\algo{chal} = \hash_p(\inst, a, b)$, and
accepts if and only if the following holds:
%
\begin{align*}
%
  a x^{\algo{chal}} &= u^\eta w^\nu \ ; \ b \hgen^{-\algo{chal}} = w^\xi
q_\ell^\pi \prod_{i=0}^{\ell-1} q_i^{\pi_i}.
  %
\end{align*}

\paragraph{Splitting the proof between two provers.} We now instantiate
$\algo{JProve}_0$ and $\algo{JProve}_1$, following the approach introduced
in~\cite{USENIX:HesSinSor23}. Here $\algo{Prv}_0$ is a lightweight device that
knows $\msg_{\ell-1}$, and $\algo{Prv}_1$ is computationally more powerful and
knows $(\varepsilon, s, \msg_0, \ldots, \msg_{\ell-2})$. Rather than proving
relation~\eqref{eq:BBSrel}, $\algo{Prv}_0$ and $\algo{Prv}_1$ jointly prove the
relation as follows:
%
\begin{align*}
%
  \com &= q_{\ell-1}^{\msg_{\ell-1}} q_\ell^\rho \ \wedge \ \pairing(u,
\tilde{t}) \pairing(x, \ggen) = 1 \ \wedge \\
  %
  x &= u^\varepsilon w^\iota \ \wedge \ (\com \hgen)^{-1} = w^\lambda
q_\ell^{s'} \prod_{i=0}^{\ell-2} q_i^{\msg_i}.
  %
\end{align*}
%
Here, $\inst = (\pk_\BBS, u, w, x, \com)$ and $\witness = (\varepsilon, s',
\iota, \lambda, \msg, \rho)$. The protocol is given in
Primitive~\ref{prim:bbs-joint-proof}. All $\algo{Prv}_1$ receives from
$\algo{Prv}_0$ is the commitment $\com$, its randomness $\rho$, and a proof of
knowledge of the opening. Since $\rho$ is disclosed, $\algo{Prv}_1$ holds
$q_{\ell-1}^{\msg_{\ell-1}}$.  As noted in~\Cref{sec:joint-pok-def},
$\msg_{\ell-1}$ is therefore protected from the more powerful prover by the
hardness of computing discrete logarithms in $\group{1}$.

Primitive~\ref{prim:bbs-joint-proof} is the construction
of~\cite{USENIX:HesSinSor23}, and satisfies the completeness, soundness, and
zero knowledge properties of~\Cref{sec:joint-pok-def} as shown there. We use it
with one simplification that our setting permits. In the original,
$\algo{Prv}_0$ derives $\rho$ by applying a pseudorandom function to a nonce,
under a key shared between the two provers. That shared key resolves a tension
we do not face. Their $\com$ must be randomised at every presentation, since an
unrandomised commitment to $\msg_{\ell-1}$ is a static identifier for
$\algo{Prv}_0$. Yet $\algo{Prv}_1$ must still recognise which of its stored
witnesses the incoming $\com$ belongs to, without the two provers ever
communicating directly. Our two provers are paired at registration and do
communicate directly, so $\algo{Prv}_0$ samples $\rho$ afresh and sends it. The
original construction would serve here unchanged.

We apply the primitive to a relation carrying one further conjunct, the opening
of the payer's pseudonym $c$. This is a conjunction of sigma protocols rather
than a change to the primitive. The conjunct is discharged by $\algo{Prv}_1$
alone and concerns no part of $\witness_0$.

\begin{primitivebox}{Joint proof of knowledge of a BBS+
signature}{bbs-joint-proof}

  \begin{description}

    \item[$\Pi \leftarrow \algo{Prv}_0(\msg_{\ell-1}, \rho) \leftrightarrow
\algo{Prv}_1(\varepsilon, s, \msg_0, \ldots, \msg_{\ell-2})$ :] The two provers
jointly compute $\Pi$ as follows.

      \begin{enumerate}

	\item $\algo{Prv}_0$ computes $\com = q_{\ell-1}^{\msg_{\ell-1}}
q_\ell^\rho$ and produces a proof $\Pi_0$ that $\inst_0 = (q_{\ell-1}, q_\ell,
\com)$ and $\witness_0 = (\msg_{\ell-1}, \rho)$ satisfy $\com =
q_{\ell-1}^{\msg_{\ell-1}} q_\ell^\rho$: it selects random $(\gamma_{\ell-1},
\upsilon)$, computes $c = q_{\ell-1}^{\gamma_{\ell-1}} q_\ell^\upsilon$,
$\algo{chal}_0 = \hash_p(\inst_0, c)$, $\pi_{\ell-1} = \gamma_{\ell-1} +
\algo{chal}_0 \msg_{\ell-1}$, and $\zeta = \upsilon + \algo{chal}_0 \rho$, and
sends $\algo{Prv}_1$ the tuple $(\com, \rho, \Pi_0)$ with $\Pi_0 = (c,
\pi_{\ell-1}, \zeta)$.

	\item $\algo{Prv}_1$ computes a proof $\Pi_1$ that, for $\inst_1 =
(\pk_\BBS, u, w, x, \com)$ and $\witness_1 = (\varepsilon, s', \iota, \lambda,
\msg[{:}\ell-1])$ with $s' = s - \rho$, the following holds:
          %
	  \begin{equation} \label{eq:reljointproof}
            %
	    x = u^\varepsilon w^\iota \ \wedge \ (\com \hgen)^{-1} = w^\lambda
\prod_{i=0}^{\ell-2} q_i^{\msg_i} q_\ell^{s'}.
            %
	  \end{equation}
          %
	  To do so, it selects random $(\alpha, \beta)$ and $(\delta, \gamma_0,
\ldots, \gamma_{\ell-2}, \gamma)$, computes $a = u^\alpha w^\beta$ and $b =
w^\delta q_\ell^\gamma \prod_{i=0}^{\ell-2} q_i^{\gamma_i}$, and derives the
following:
          %
	  \begin{align*} \algo{chal}_1 &= \hash_p(\inst_1, a, b) \ ; \ \eta =
\alpha + \algo{chal}_1 \varepsilon \ ; \ \nu = \beta + \algo{chal}_1 \iota \ ; \
\xi = \delta + \algo{chal}_1 \lambda \ ; \\
            %
	    \pi_i &= \gamma_i + \algo{chal}_1 \msg_i, \ \forall i \in [\ell-1] \
; \ \pi = \gamma + \algo{chal}_1 s'.
            %
	  \end{align*}

	\item $\algo{Prv}_1$ outputs $\Pi = (\com, a, b, c, \eta, \nu, \xi,
\pi_0, \ldots, \pi_{\ell-1}, \pi, \zeta)$.

      \end{enumerate}

    \item[$\{0,1\} \leftarrow \algo{JVerify}(\inst, \Pi)$ :] On input $\inst =
(\pk_\BBS, u, w, x, \com)$ and $\Pi$, do:

      \begin{enumerate}

	\item Recompute $\algo{chal}_0 = \hash_p((q_{\ell-1}, q_\ell, \com), c)$
and $\algo{chal}_1 = \hash_p(\inst, a, b)$. The two sub-proofs carry separate
challenges, since $\algo{Prv}_0$ commits before $\algo{Prv}_1$ has formed
$\inst$.

	\item Return the conjunction of the following three checks:
          %
	  \begin{align*}
          %
	    &\pairing(u, \tilde{t}) \pairing(x, \ggen) \checkcheck 1 \ ; \\
            %
	    &c \,\com^{\algo{chal}_0} \checkcheck q_{\ell-1}^{\pi_{\ell-1}}
q_\ell^{\zeta} \ ; \\
            %
	    &a x^{\algo{chal}_1} \checkcheck u^\eta w^\nu \ \wedge \ b (\com
\hgen)^{-\algo{chal}_1} \checkcheck w^\xi q_\ell^\pi \prod_{i=0}^{\ell-2}
q_i^{\pi_i}.
            %
	  \end{align*}
          %
	  The third check is relation~\eqref{eq:reljointproof} in verified form.

      \end{enumerate}

  \end{description}

\end{primitivebox}

\paragraph{Mapping to our solution.} In our protocol, the user's secure element
and mobile device jointly produce a proof of knowledge of the user's credential.
The witness is split so that the secure element contributes only knowledge of
the user's secret key, and the device supplies the rest. Concretely, they
jointly produce $\sigma \leftarrow \NIZKprove(\crs, \inst, \witness)$, where
$\inst = (\gen_0, \gen_1, \pk_\BBS, u, w, x, c)$ with $\pk_\BBS = (\ggen,
\tilde{t}, q_0, q_1, q_2, \hgen)$, $\witness = (\varepsilon, \iota, \lambda, s,
\algo{uid}, \sk, r)$, and $\rel(\inst, \witness) = 1$ is equivalent to the
following:
%
\begin{align*}
%
  &\pairing(u, \tilde{t}) \pairing(x, \ggen) = 1 \ \wedge \ x = u^\varepsilon
w^\iota \\ &\wedge \ \hgen^{-1} = w^\lambda q_0^{\algo{uid}} q_1^\sk q_2^s \
\wedge \ c = \gen_0^{\algo{uid}} \gen_1^r.
  %
\end{align*}
%
That is, this is an instantiation of relation~\eqref{eq:BBSrel} with $\msg =
(\algo{uid}, \sk)$ and $c = \gen_0^{\algo{uid}} \gen_1^r$. In this
instantiation, $\algo{se}$ executes $\algo{Prv}_0$, outputting $\com = q_1^\sk
q_2^\rho$, $\rho$, and proof $\Pi_0$ computed as above except that the challenge
$\algo{chal}_0$ also hashes the message $m_i$ of Protocol~\ref{prot:offline}
that the signature of knowledge is taken over.  This is what binds the payee's
pseudonym to the proof, and it is the only part of the message that the secure
element needs to see.  $\algo{dev}$ computes a proof $\Pi_1 \leftarrow
\NIZKprove(\crs, \inst_1, \witness_1)$, for $\inst_1 = (\gen_0, \gen_1, \pk, u,
w, x, \com, c)$, $\witness_1 = (\varepsilon, \iota, \lambda, \algo{uid}, s',
r)$, and relation as follows:
%
\begin{equation*}
%
  x = u^\varepsilon w^\iota \ \wedge \ (\com \hgen)^{-1} = w^\lambda
q_0^{\algo{uid}} q_2^{s'} \ \wedge \ c = \gen_0^{\algo{uid}} \gen_1^r.
  %
\end{equation*}
%

\paragraph{Merging the audit proof.} The audit proof $\Gamma_i$ of
Protocol~\ref{prot:offline} proves that $\ciph_i$ encrypts the value committed
in $c$, on witness $(\algo{uid}, r, \rho_i)$. Its first conjunct is the opening
of $c$, which is already the third conjunct of $\Pi_1$, and both are computed by
$\algo{dev}$ on the same witness. We therefore compute the two as one sigma
protocol under a single challenge $\algo{chal}_1$, adding only the two
ciphertext conjuncts and the response $\pi_{\rho}$ for the encryption
randomness. This saves one opening of $c$ per hop, in both proving and
verification, relative to computing $\Gamma_i$ and $\sigma_i$ independently.

\subsection{Instantiating the token signatures}

The central bank signs each token at withdrawal, and must do so without learning
what it is signing. We use the Pointcheval-Sanders signature scheme, which also
signs vectors of messages, and extend it to sign committed inputs.  The
algorithms are given in
Primitive~\ref{prim:ps}~\cite{RSA:PoiSan16,RSA:PoiSan18}.

\begin{primitivebox}{The Pointcheval-Sanders signature scheme}{ps}

  \begin{description}

    \item[$(\sk, \pk) \leftarrow \PSgen(\secparam, \ell)$ :] On input security
parameter $\secparam$ and integer $\ell$, do the following:

      \begin{enumerate}

	\item Sample $(\chi, \lambda_0, \ldots, \lambda_{\ell-1}, \iota) \sample
\Zp^{\ell+2}$.

	\item Compute $\tilde{x} = \ggen^\chi$, $\tilde{y}_i =
\ggen^{\lambda_i}$ for all $i \in [\ell]$, and $\tilde{z} = \ggen^\iota$.

	\item Return $\sk = (\chi, \vec{\lambda}, \iota)$ and $\pk = (\tilde{x},
\vec{\tilde{y}}, \tilde{z})$.

      \end{enumerate}

    \item[$\psi \leftarrow \PSsign(\sk, \msg)$ :] On input secret key $\sk$ and
a vector of $\ell$ messages $\msg$, do:

      \begin{enumerate}

	\item Sample $k \sample \group{1}$ and $\omega \sample \Zp$.

	\item Compute $l = k^{\chi + \vec{\lambda} \cdot \msg + \iota\omega}$.

	\item Return $\psi = (\omega, k, l)$.

      \end{enumerate}

    \item[$\{0,1\} \leftarrow \PSverify(\pk, \msg, \psi)$ :] On input public key
$\pk$, vector of messages $\msg$, and signature $\psi$, return the result of the
following:
      %
      \begin{equation*}
      %
	\pairing(k, \tilde{x} \vec{\tilde{y}}^\msg \tilde{z}^\omega) \checkcheck
\pairing(l, \ggen).
            %
      \end{equation*}
      %

  \end{description}

\end{primitivebox}

The resulting signatures are existentially unforgeable under the $q$-modified
strong Diffie-Hellman ($q$-MSDH) assumption of
Assumption~\ref{asm:qmsdh}~\cite{RSA:PoiSan18}.

\begin{assumptionbox}{$q$-Modified Strong Diffie-Hellman ($q$-MSDH)}{qmsdh}

  The $q$-MSDH assumption holds if, given $(\gen, \gen^\chi, \ggen, \ggen^\chi,
\ggen^{\omega_1}, \ldots, \ggen^{\omega_q}, \gen^{1/(\chi+\omega_1)}, \ldots,
\gen^{1/(\chi+\omega_q)})$ for uniformly random $\chi, \omega_1, \ldots,
\omega_q \sample \Zp$, no PPT adversary can output a tuple $(\omega, \gen^\nu,
\gen^{\nu/(\chi+\omega)})$ for $\nu \neq 0$ and $\omega \notin \{\omega_1,
\ldots, \omega_q\}$ with non-negligible probability.

\end{assumptionbox}

Pointcheval-Sanders signatures can also be extended to support blind signing: a
signer can sign a committed vector without learning any information about
it~\cite{NDSS:SABMD19}.

\paragraph{Blind signing over committed inputs.} Let $\vec{\gen} = (\gen_0,
\ldots, \gen_\ell) \in \group{1}^{\ell+1}$ be a vector of $\ell+1$ random
generators and $\msg = (\msg_0, \ldots, \msg_{\ell-1}) \in \Zp^\ell$ a vector of
$\ell$ messages. Signatures over committed inputs are produced via an
interactive protocol between a signer and a
recipient~\cite{SCN:CamLys02,ASIACCS:LMPY16}. The signer signs a vector of
committed messages, and only the recipient learns their content. The ciphertexts
below are written $\hat{\ciph}$ throughout, and are encryptions under a key of
the recipient's own, generated freshly for each execution so that two runs
cannot be linked by their public key. They are unrelated to the audit
ciphertexts $\ciph_i$ of Protocol~\ref{prot:offline}, which are under the
auditor's key.

The recipient selects an ElGamal secret key $\mu \sample \Zp$ with public key
$(\gen, b = \gen^\mu)$, computes $\com = \vec{\gen}^{\msg\|r}$, $k =
\hash_{\group{1}}(\com)$, and $\ell$ ciphertexts $\hat{\ciph}_i =
(\gen^{\rho_i}, k^{\msg_i} b^{\rho_i})$ for $i \in [\ell]$, together with a
proof $\Pi \leftarrow \NIZKprove(\crs, \inst, \witness)$ for $\inst =
(\vec{\gen}, \gen, b, \com, k, (\hat{\ciph}_i)_{i \in [\ell]})$, $\witness =
((\msg_i, \rho_i)_{i \in [\ell]}, r)$, and relation as follows:
%
\begin{align*}
%
  \com = \vec{\gen}^{\msg\|r} \ \wedge \ \hat{\ciph}_i = (\gen^{\rho_i},
k^{\msg_i}b^{\rho_i}), \ \forall i \in [\ell].
  %
\end{align*}
%
It then sends the signer $(\com, \hat{\ciph}_0, \ldots, \hat{\ciph}_{\ell-1},
\gen, b, \Pi)$, and no other value: neither $\msg$ nor $r$ nor any $\rho_i$
leaves the recipient. The signer recomputes $k = \hash_{\group{1}}(\com)$,
verifies $\Pi$ against $\com$, $k$, and $(\gen, b)$, and rejects if it does not
hold. Otherwise, it parses each $\hat{\ciph}_i$ as $(e_{i,1}, e_{i,2})$,
randomly selects $\omega \in \Zp$, and computes the following:
%
\begin{align*}
%
  \hat{\ciph} &= (e_1, e_2) = \Bigl(\prod_{i=0}^{\ell-1} e_{i,1}^{\lambda_i}, \
k^\chi \prod_{i=0}^{\ell-1} e_{i,2}^{\lambda_i} k^{\omega\iota}\Bigr).
  %
\end{align*}
%
It then sends $(\omega, \hat{\ciph})$ to the recipient. The recipient decrypts
$\hat{\ciph} = (e_1, e_2)$ by computing $l = e_2 / e_1^\mu$, and outputs $\psi =
(\omega, k, l)$ if it is a valid Pointcheval-Sanders signature on $\msg$. If the
signer followed the protocol honestly, $l = k^\chi \bigl(\prod_{i=0}^{\ell-1}
k^{\lambda_i \msg_i}\bigr) k^{\omega\iota}$. Note that the blind variant fixes
$k = \hash_{\group{1}}(\com)$ rather than sampling it as $\PSsign$ does, since
the signer never sees $\msg$ and both parties must agree on $k$. Unforgeability
is preserved, as $k$ is still a group element neither party
controls~\cite{NDSS:SABMD19}.

This protocol adapts readily to partially opened vectors: for $\algo{Disc}
\subseteq [\ell]$ with $(\msg_i)_{i \in \algo{Disc}}$ revealed to the signer,
the recipient computes $\com = \prod_{i \in [\ell]\setminus\algo{Disc}}
\gen_i^{\msg_i}$ and a proof for $\inst = (\algo{Disc}, \vec{\gen}, \gen, b,
\com, k, (\hat{\ciph}_i)_{i \in [\ell]\setminus\algo{Disc}})$, $\witness =
((\msg_i, \rho_i)_{i \in [\ell]\setminus\algo{Disc}}, r)$, and relation as
follows:
%
\begin{align*}
%
  \com &= \gen_\ell^r \prod_{i \in [\ell]\setminus\algo{Disc}} \gen_i^{\msg_i} \
\wedge \ \hat{\ciph}_i = (\gen^{\rho_i}, k^{\msg_i} b^{\rho_i}), \ \forall i \in
[\ell]\setminus\algo{Disc}.
  %
\end{align*}
%
It sends $\msg_i$ for $i \in \algo{Disc}$ in the clear, and the signer computes
the following:
%
\begin{align*}
%
  \hat{\ciph} &= (e_1, e_2) = \Bigl(\prod_{i \in [\ell]\setminus\algo{Disc}}
e_{i,1}^{\lambda_i}, \ k^\chi \prod_{i \in \algo{Disc}} k^{\msg_i \lambda_i}
\prod_{i \in [\ell]\setminus\algo{Disc}} e_{i,2}^{\lambda_i}
k^{\omega\iota}\Bigr).
  %
\end{align*}
%
This protocol adapts directly to the withdrawal step of
Protocol~\ref{prot:offline}: we execute it on the vector of messages
$(\algo{uid}_0, v, s_0)$ with $\algo{Disc} = \{1\}$ and randomness $r = 0$, so
that the request commitment $C$ of Protocol~\ref{prot:offline} is $\com$ and the
two ciphertexts sent are $\hat{\ciph}_0$ and $\hat{\ciph}_2$.

\paragraph{Proof of withdrawal.} Let $\tau = (v, c = \gen_0^{\algo{uid}}
\gen_1^s)$ be a token, and let $\pk_{\algo{PS}} = (\tilde{x}, \tilde{y}_0,
\tilde{y}_1, \tilde{y}_2, \tilde{z})$ be the public key of the central bank
$\algo{cb}$. Let $\psi = (\omega, k, l)$ be a Pointcheval-Sanders signature that
verifies $1 \gets \PSverify(\pk_{\algo{PS}}, (\algo{uid}, v, s), \psi)$, so
$\pairing(k, \tilde{x} \tilde{y}_0^{\algo{uid}} \tilde{y}_1^v \tilde{y}_2^s
\tilde{z}^\omega) = \pairing(l, \ggen)$. To prove that $\tau$ was issued by
$\algo{cb}$, the token's owner uses their mobile device $\algo{dev}$ to compute
$\Sigma$ as follows.

\begin{protocolbox}{Proof of withdrawal}{ps-proof}

  \begin{description}

    \item[$\Sigma \leftarrow \algo{ProveWithdrawal}(\pk_{\algo{PS}}, \tau,
\omega, \algo{uid}, s)$ :] On input the central bank's public key
$\pk_{\algo{PS}}$, token $\tau = (v, c)$, and witness $(\omega, \algo{uid}, s)$
for a Pointcheval-Sanders signature $\psi = (\omega, k, l)$ on $(\algo{uid}, v,
s)$, do:

      \begin{enumerate}

	\item Select random $\iota \in \Zp$ and compute $(k^*, l^*) = (k^\iota,
l^\iota)$.

	\item Generate $\Delta \leftarrow \NIZKprove(\crs, \inst, \witness)$,
with $\inst = (\gen_0, \gen_1, \tilde{x}, \tilde{y}_0, \tilde{y}_1, \tilde{y}_2,
\tilde{z}, v, c, k^*, l^*)$, $\witness = (\omega, \algo{uid}, s)$, and relation
as follows:
          %
	  \begin{align*}
          %
	    \pairing(l^*, \ggen) &= \pairing(k^*, \tilde{x}
\tilde{y}_0^{\algo{uid}} \tilde{y}_1^v \tilde{y}_2^s \tilde{z}^\omega) \ \wedge
\ c = \gen_0^{\algo{uid}} \gen_1^s.
            %
	  \end{align*}
          %

	\item Return $\Sigma = (k^*, l^*, \Delta)$.

      \end{enumerate}

  \end{description}

\end{protocolbox}

Since $s_0$ is random and not known, $(k^*, l^*)$ leaks no information about
$\algo{uid}_0$. To compute $\Delta$, $\algo{dev}$ first computes $\tilde{c} =
\tilde{y}_0^{\algo{uid}} \tilde{y}_2^s \tilde{z}^\omega$. It randomly selects
$\alpha, \beta, \delta \in \Zp$ and computes $a = \gen_0^\alpha \gen_1^\beta$,
$\tilde{a} = \tilde{y}_0^\alpha \tilde{y}_2^\beta \tilde{z}^\delta$, and
challenge $\gamma = \hash_p(\inst', \tilde{c}, a, \tilde{a})$. It then computes
$\pi_{\algo{uid}} = \alpha + \gamma \algo{uid}$, $\pi_s = \beta + \gamma s$, and
$\pi_\omega = \delta + \gamma \omega$, and sets $\Delta = (\tilde{c}, a,
\tilde{a}, \pi_{\algo{uid}}, \pi_s, \pi_\omega)$. These are standard Schnorr
proofs for $\inst' = (\gen_0, \gen_1, \tilde{y}_0, \tilde{y}_2, \tilde{z}, c,
\tilde{c})$, $\witness' = (\omega, \algo{uid}, s)$, and $\rel(\inst', \witness')
= 1$ equivalent to $c = \gen_0^{\algo{uid}} \gen_1^s \wedge \tilde{c} =
\tilde{y}_0^{\algo{uid}} \tilde{y}_2^s \tilde{z}^\omega$.

If $(\omega, k^*, l^*)$ is a valid Pointcheval-Sanders signature on
$(\algo{uid}, v, s)$ under $\pk_{\algo{PS}}$, then $\pairing(l^*, \ggen) =
\pairing(k^*, \tilde{x} \tilde{c} \tilde{y}_1^v)$. A verifier is given $\Sigma =
(k^*, l^*, \Delta)$ and token $\tau = (v, c)$. It checks that $\tau$ was issued
by $\algo{cb}$ by parsing $\Delta = (\tilde{c}, a, \tilde{a}, \pi_{\algo{uid}},
\pi_s, \pi_\omega)$, computing $\gamma = \hash_p(\inst', \tilde{c}, a,
\tilde{a})$, and verifying the following:
%
\begin{align*}
%
  \pairing(l^*, \ggen) &\checkcheck \pairing(k^*, \tilde{x} \tilde{c}
\tilde{y}_1^v) \ \wedge \ c^\gamma a \checkcheck \gen_0^{\pi_{\algo{uid}}}
\gen_1^{\pi_s} \ \wedge \ \tilde{c}^\gamma \tilde{a} \checkcheck
\tilde{y}_0^{\pi_{\algo{uid}}} \tilde{y}_2^{\pi_s} \tilde{z}^{\pi_\omega}.
  %
\end{align*}

\subsection{The instantiated protocol}\label{sec:offline-full-inst}

The pieces above instantiate every generic component of
Protocol~\ref{prot:offline}, as summarised in~\Cref{tab:offline-instantiation}.

\begin{table}[t] \centering \small \begin{tabular}{l|l|l} \toprule Component &
Instantiation & Assumption \\ \midrule Commitment scheme & Pedersen in
$\group{1}$ & DL in $\group{1}$ \\ Credential signature & BBS+ & $q$-SDH \\
Blind signature & Pointcheval-Sanders & $q$-MSDH \\ Public key encryption &
ElGamal in $\group{1}$ & DDH in $\group{1}$ \\ Joint proof of knowledge &
Primitive~\ref{prim:bbs-joint-proof} & ROM \\ Proof of withdrawal &
Protocol~\ref{prot:ps-proof} & ROM \\ \bottomrule \end{tabular}
\caption{Instantiation of each generic component of Protocol~\ref{prot:offline},
with the assumption under which it is secure.  All proofs are sigma protocols
compiled with Fiat-Shamir, so each carries the random oracle model in addition
to the discrete logarithm assumption in $\group{1}$.}
\label{tab:offline-instantiation} \end{table}

A user's pseudonym is the Pedersen commitment $\com = \gen_0^{\algo{uid}}
\gen_1^s$, which is perfectly hiding and computationally binding under the
discrete logarithm assumption in $\group{1}$. The registration authority signs
the pair $(\algo{uid}, \sk_{\algo{usr}})$ with BBS+, blindly, so that
$\sk_{\algo{usr}}$ never leaves the secure element. The central bank signs the
token vector $(\algo{uid}_0, v, s_0)$ with Pointcheval-Sanders over committed
inputs, learning neither $\algo{uid}_0$ nor $s_0$. The audit ciphertext is an
ElGamal encryption of the payer's identifier under $\pk_\auditor$, which admits
an efficient sigma protocol proving that it encrypts the value committed in the
payer's pseudonym. That is the proof $\Gamma_i$, merged with $\sigma_i$ as
described above. The signature of knowledge $\sigma_i$ that authorises each
payment and each deposit is Primitive~\ref{prim:bbs-joint-proof}, and the proof
of withdrawal $\Sigma$ is produced by the device alone using
Protocol~\ref{prot:ps-proof}.

Every proof in the instantiation is a sigma protocol over discrete logarithms
with the Fiat-Shamir heuristic applied, which is simulation extractable in the
random oracle model~\cite{C:FiaSha86,INDOCRYPT:FKMV12}. This includes
Primitive~\ref{prim:bbs-joint-proof}, which is a conjunction of such protocols,
so the same argument gives simulation extractability rather than only the
soundness established in~\cite{USENIX:HesSinSor23}. This gives the following.

\begin{theorem}\label{thm:offline-instantiation-security}
%
  If the discrete logarithm and DDH assumptions hold in $\group{1}$, the $q$-SDH
assumption holds for the BBS+ signature scheme, and the $q$-MSDH assumption
holds for the Pointcheval-Sanders signature scheme, then the instantiated
offline payment protocol securely realises ideal functionality
$\idfu{\algo{op}}$ in the random oracle model.
%
\end{theorem}
